% Paper 11 v2 (Route A): Silent-Failure Taxonomy and Audit Protocol
% Target: COMPEL (Elsevier, elsarticle one-column). Two-column IEEE draft kept in main.tex
\documentclass[preprint,12pt]{elsarticle}
\usepackage{amsmath,graphicx,booktabs,url}
\journal{COMPEL - The international journal for computation and
mathematics in electrical and electronic engineering}
\begin{document}
\begin{frontmatter}
\title{Silent Failures in Frequency-Domain Finite-Element Analysis of
Cage Induction Motors: A Taxonomy, Forensic Case Studies, and an Audit
Protocol}
\author[a]{Anonymous\corref{cor1}}
\ead{author@example.org}
\cortext[cor1]{Corresponding author.}
\affiliation[a]{organization={ChinaSim Stdio Motor Workbench},
addressline={In-house research}, city={Beijing}, country={China}}
\begin{abstract}
Finite-element analysis of cage induction motors couples magnetic field
unknowns with circuit unknowns and, in multi-pole geometries, periodic
sector constraints. We show that this class of simulation fails in a
characteristic way: the solver converges, network conservation residuals
are exactly zero, and no warning is emitted --- yet the reported quantity
is wrong. We document three forensic cases on an open-source GetDP~4.0.0
toolchain: (A)~a sparse direct solver (MUMPS) aborting on a workspace-budget
limit (\texttt{INFO(1)=-19}, ICNTL(23)) on a system that an independent
factorization solves in 26\,s (residual $2.28\times10^{-5}$) --- the
abort is silently propagated downstream as an all-NaN field; (B)~a ``half-field''
solution at an intermediate geometric scale, physically wrong by 42\% with
no numerical diagnostic; (C)~sector reduction violating a uniform-sign
condition on harmonic commensurability, inflating the air-gap fundamental
by $1.44\times$ and torque by $1.54\times$ while KCL remains $0.000$. From the taxonomy we derive a
three-check audit protocol --- cross-solver verdict, scale
self-consistency band, and sector feasibility --- each with an executable
criterion, and we release reference implementations. We delimit what is
folklore (sign intuition for periodic boundaries) from what is not (the
silence of the failures, their quantitative extent, and the audit
coverage gap). All quantitative claims trace to archived raw data.
\end{abstract}
\begin{keyword}
induction motor \sep finite element \sep silent failure \sep direct
solver \sep periodic boundary condition \sep verification protocol
\end{keyword}
\end{frontmatter}

\section{Introduction}
Application-level silent failures --- wrong results produced by
successfully completed solves --- are studied extensively at the hardware
level as silent data corruption in high-performance computing
\cite{sdc2011,sdc2014}, and at the mathematical level as result
verification of differential-equation solvers \cite{resver2004}. Between
these two layers sits a third failure class that, to our knowledge, has
not been systematized: \emph{modeling- and configuration-level} silent
failures of engineering simulation chains, where the linear algebra is
performed correctly on the wrong system, or the right system under an
infeasible constraint. During a frequency-domain cage-motor study on an
open-source toolchain we collected three such failures with complete
forensics. This paper (i)~organizes them into a taxonomy
(Sec.~\ref{sec:tax}), (ii)~presents the forensics with all numbers
traceable to archived raw data (Secs.~\ref{sec:caseA}--\ref{sec:caseC}),
(iii)~distills a three-check audit protocol with executable criteria
(Sec.~\ref{sec:audit}), and (iv)~delimits the folklore boundary
(Sec.~\ref{sec:folk}).

\section{Failure Taxonomy}\label{sec:tax}
\begin{table}[t]\centering\caption{Taxonomy of application-level silent
failures observed in the study.}\label{tab:tax}
\begin{tabular}{p{2.6cm}p{3.2cm}p{3.4cm}p{2.6cm}}\toprule
class & mechanism & detectable by & forensic case \\\midrule
F1 false infeasibility & solver-specific factorization failure on a
workspace-budget abort, silently propagated & cross-solver verdict & A
(Sec.~\ref{sec:caseA})\\
F2 silent degradation & scale-dependent conditioning yields a
pathological solution without NaNs & scale self-consistency band & B
(Sec.~\ref{sec:caseB})\\
F3 infeasible constraint & periodic constraint with conflicting
harmonic signs projects onto a nonphysical subspace & sector feasibility
(uniform-sign) test & C (Sec.~\ref{sec:caseC})\\\bottomrule
\end{tabular}\end{table}
Table~\ref{tab:tax} defines the three classes. Their common signature is
the \emph{absence of alarms}: residual checks, conservation laws, and
solver exit codes all pass; only cross-validation of physical quantities
reveals the failure.

\begin{figure}[t]\centering\includegraphics[width=0.85\columnwidth]{../../figure/fig1_threshold_ladder.pdf}
\caption{Failure threshold ladder for the scaled cage model: MUMPS outcome vs.\ geometry scale factor and solver. [Data: data/raw/case1\_scale\_ladder.]\label{fig:threshold}}\end{figure}
\section{Case A: False Infeasibility}\label{sec:caseA}
A scaled 16-bar cage model ($\chi$ = geometric scale factor on the
12-slot/4-pole ASM baseline) is solved frequency-domain at slip $s=1$.
For $\chi\le1.75$ the MUMPS direct solve succeeds; at $\chi=2$ it
terminates with error $-19$ --- per the MUMPS 5.8.2 sources
(\texttt{dfac\_b.F:283}) a workspace-budget abort (ICNTL(23) exceeded), not a singularity diagnostic --- and the all-NaN field is silently
propagated through the pipeline in 37\,s; and at $\chi=3.067$ the same failure persists under
out-of-core settings (81.2\,s, all-NaN). Four exclusions pin the cause:
(i) a pure-field bisection with all circuit terms commented out still
reports $-19$; (ii) mesh quality, group integrity, and Dirichlet
placement are identical across the threshold; (iii) bar-ratio
normalization does not rescue $\chi=3.067$; and (iv) the serial
SuperLU LU of the identical matrix converges in 26\,s to residual
$2.28\times10^{-5}$, with torque $T=0.7477$\,N\,m and $B_{g1}=
0.0158$\,T consistent with the $\chi=1$ baseline under the gap scaling
of the family. The failure is therefore solver-specific
(infeasibility is false), and a later confound refinement showed the
historical $\chi=2$ failure also carried a bar-ratio distortion:
with the ratio normalized, $\chi=2$ solves in both solvers ($T=
0.2495$\,N\,m, $B_{g1}=0.0363$\,T), moving the failure band to
$(\chi\,2.0, 3.067]$. All numbers trace to data/raw/case1\_*.]
\begin{figure}[t]\centering\includegraphics[width=0.9\columnwidth]{../../figure/fig2_cross_solver.pdf}
\caption{Cross-solver decision: same matrix aborted by MUMPS on ICNTL(23) budget, converged by SuperLU, quantified.
[Data: data/raw/case1\_cross\_solver.]\label{fig:cross}}\end{figure}

\section{Case B: Silent Degradation}\label{sec:caseB}
At $\chi=1.5$ the solve completes without NaNs but $B_{g1}=0.0663$\,T is
$58\%$ of the self-consistent value ($0.1141$ at $\chi=1$ and $0.1161$
at $\chi=1.75$; both at slip $s=0.05$, \(j_s=3\,\mathrm{A/mm^2}\)): a $42\%$
silent error bracketed by two healthy runs. (All Case-A solves at $s=1$
use the separate baseline $B_{g1}=0.0579$\,T; the two baselines refer to
different operating points of the same machine.)
[Data: data/raw/case1\_scale\_ladder.]
\begin{figure}[t]\centering\includegraphics[width=0.9\columnwidth]{../../figure/fig3_sector_contamination.pdf}
\caption{Sector contamination: sectorized vs.\ full-model spectra, $\times1.44$ (flux) / $\times1.54$ (torque) shifts with zero KCL residual. [Data: data/raw/case2\_sector.]\label{fig:sector}}\end{figure}
\section{Case C: Infeasible Constraint}\label{sec:caseC}
A $90^\circ$ sector of the same machine, linked by a single periodic
constraint, passes every solver-level check (SuperLU converges in
6.2\,s, KCL residual $0.000$) yet reports $T=0.0605$\,N\,m and
$B_{g1}=0.0833$\,T against the full-model $T=0.0393$\,N\,m,
$B_{g1}=0.0579$\,T: a $\times1.54$ (torque, $0.0605/0.0393$) and $\times1.44$
(flux density, $0.0833/0.0579$) (flux-density) signed shift with no alarm. The mechanism is order
incompatibility: the fundamental (order 2) is anti-periodic over
$90^\circ$, whereas 12k-order stator-slot harmonics (k odd) are
integer-periodic, so one Link sign cannot serve both --- the sector
constraint is structurally infeasible for the coupled cage model, not
merely inaccurate.
[uniform-sign proposition + $90^\circ$/12s4p sign table Fig.4 + sector
vs.\ full-model numbers $1.44/1.54\times$ + KCL $=0.000$ ---
Data: data/raw/case2\_sector.]
\section{Audit Protocol}
The protocol ships as an executable reference implementation
(\texttt{src/audit\_protocol.py}, bundled): steps A1--A5 consume the
solver reports directly and emit a verdict with per-step evidence; its
doctest suite is driven by the archived FAIL-branch and PASS-branch
reports of Section~\ref{sec:audit} (nine assertions, all green), so the
auditor itself is regression-tested against the same forensic data.\label{sec:audit}
Check~1 (cross-solver): on any direct-solver failure, re-solve with an
independent factorization and compare physical quantities, not only
residuals. Check~2 (scale band): in any geometry/scale sweep, flag runs
whose key quantity deviates $>10\%$ from the interpolation of healthy
neighbors. Check~3 (sector feasibility): before sector reduction, verify
$m\,\theta_s \equiv 0$ or $180^\circ \pmod{360^\circ}$ with the
\emph{same} sign for the fundamental and all slot/teeth/cage orders;
reject infeasible sectors. Reference implementations ship with the
paper sources.
\section{What Is Obvious and What Is Not}\label{sec:folk}
That periodic boundaries should match harmonic periods is folklore, and
simulator guides warn about it. What is not folklore, and what motivates
this paper: (i)~the failure is \emph{silent} --- every standard check
passes; (ii)~the error is \emph{large and signed} ($+44\%$ on the
fundamental); (iii)~in coupled cage models the infeasibility can be
\emph{structural} (single Link sign cannot serve fundamental and teeth
harmonics simultaneously), so no sector angle within a broad family is
admissible;
\begin{figure}[t]\centering\includegraphics[width=0.9\columnwidth]{../../figure/fig4_sign_table.pdf}
\caption{Uniform-sign feasibility table: single Link sign vs.\ significant mode order for the 12-slot/4-pole family.
[Data: data/raw/case2\_sector.]\label{fig:sign}}\end{figure}
 (iv)~an analogous silence holds for F1 and F2, and no
published protocol covers the three jointly for electrical-machine FE.
\section{Implications for Simulation and Geometry Pipelines}
(G1) bar-ratio normalization; (G2) sector pre-check; (G3) scale-band
self-consistency; (G4) solver escalation before model changes. All four
are implemented in our toolchain and referenced from the case
forensics above.
\section{Conclusion}
Application-level silent failures in electrical-machine FE form a
coherent class with a shared signature (no alarms) and complementary
audits (cross-solver, scale band, uniform sign). The three forensic
cases, the executable protocol, and the archived data are released to
make such audits routine.
\section*{Data availability}
Raw solve logs, archived reports, generated figures, and reference
implementations are bundled with the manuscript sources
(\texttt{data/raw/}, \texttt{figure/}, \texttt{src/}).
\section*{Declaration of competing interest}
The authors declare no competing interests.
\begin{thebibliography}{9}
\bibitem{sdc2011} Detection and correction of silent data corruption for
large-scale high-performance computing, IPDPS 2011,
doi:10.1109/ipdps.2011.379.
\bibitem{sdc2014} Silent error detection in numerical time-stepping
schemes, \emph{IJHPCA}, 2014, doi:10.1177/1094342014532297.
\bibitem{resver2004} Novel approaches to numerical software with result
verification, LNCS, 2004, doi:10.1007/978-3-540-24738-8\_17\_17.
\end{thebibliography}
\end{document}
